%%%%%%%%%%%%%%%%%%%%%%%%%%%%%%%%%%%%%%%%%%%%%%%%%%%%%%%%%%%%%%%%%%%%%%
% Lớp 10 | Đại số tổ hợp | Quy tắc cộng-quy tắc nhân
% Dạng 0D8?1-2: Bài toán chỉ sử dụng quy tắc nhân
% Số câu: 64  (sinh tự động bởi tach_cau.py; từ nay thầy sửa trực tiếp file này)
%%%%%%%%%%%%%%%%%%%%%%%%%%%%%%%%%%%%%%%%%%%%%%%%%%%%%%%%%%%%%%%%%%%%%%

%%%%%%%%%%%%%%%% TRẮC NGHIỆM 4 PHƯƠNG ÁN %%%%%%%%%%%%%%%%

% nguồn: MaHoaNH_TracNghiem_Lop10
\begin{ex}%[0D8N1-2]%[ma:62761c90]
	Có bao nhiêu số tự nhiên gồm $4$ chữ số được tạo  thành từ các chữ số $1$, $2$, $3$, $4$, $5$?
	\choice 
		 {$120$}
		{$15$} 
		{\True $625$} 
		 {$32$}
	 \loigiai{
	Số các  số tự nhiên gồm $4$ chữ số được tạo thành từ các chữ số $1$, $2$, $3$, $4$, $5$ là $5^4=625$ số.}
\end{ex}

% nguồn: MaHoaNH_TracNghiem_Lop10
\begin{ex}%[0D8N1-2]%[ma:df22bf33]
	Một đội học sinh giỏi của trường THPT, gồm $5$ học sinh nam và $6$ học sinh nữ. Số cách chọn hai học sinh,  trong đó có một học sinh nam và một học sinh nữ là
	\choice 
		 {$25$}
		{$42$} 
		 {$11$}
		{\True $30$} 
	 \loigiai{
		Số cách chọn  thỏa mãn là $5\cdot 6=30$ cách.
	 }
\end{ex}

% nguồn: MaHoaNH_TracNghiem_Lop10
\begin{ex}%[0D8N1-2]%[ma:5be79610]
	Một công việc được hoàn thành bời hai hành động liên tiếp. Nếu hành động thứ nhất có $m$ cách thực hiện, hành động thứ hai có $n$ cách thực hiện (các cách thực hiện của cả $2$ hành động là khác nhau đôi một) thì số cách hoàn thành công việc  đó là
	\choice 
		{$\dfrac{m}{n}$} 
		{\True  $m\cdot n$}
		 {$m+n$}
		{$m-n$} 
	 \loigiai{
		Theo quy tắc nhân thì số cách  hoàn thành công việc đó là $m\cdot n$.
	} 
\end{ex}

% nguồn: MaHoaNH_TracNghiem_Lop10
\begin{ex}%[0D8N1-2]%[ma:6c826c69]
	Trong đội văn nghệ của nhà trường có $8$ học sinh nam và $6$ học sinh nữ. Hỏi có bao nhiêu cách chọn  ra một đôi song ca nam nữ?
	\choice 
		 {$14$}
		 {$91$}
		{\True $48$} 
		 {$182$}
	\loigiai{ 
		Theo quy tắc nhân, có $8\cdot 6=48$ cách chọn ra một đôi  song ca nam nữ.
	 }
\end{ex}

% nguồn: MaHoaNH_TracNghiem_Lop10
\begin{ex}%[0D8N1-2]%[ma:82dd1b01]
	Từ thành phố $A$ tới thành phố $B$ có $3$  con đường, từ thành phố $B$ tới thành phố $C$ có $4$ con đường. Hỏi có bao nhiêu cách đi từ $A$ tới $C$ qua $B$?
	\choice 
		 {\True $12$}
		{$24$} 
		 {$6$}
		{$7$} 
	 \loigiai{
		Từ $A$ đến $B$ có $3$ cách chọn đường  đi, từ $B$ đến $C$ có $4$ cách chọn đường đi.\\
		Vậy số cách chọn  đường đi từ $A$ đến $C$ phải đi qua $B$ là $3\cdot 4=12$ cách.
	} 
\end{ex}

% nguồn: MaHoaNH_TracNghiem_Lop10
\begin{ex}%[0D8N1-2]%[ma:ed7ca4a4]
	Một người vào cửa hàng ăn, người đó chọn thực đơn gồm $1$ món ăn trong $5$ món ăn, $1$ loại quả tráng miệng trong $4$ loại quả tráng miệng và $1$ loại nước uống trong $3$ loại nước uống. Hỏi có bao nhiêu cách chọn thực  đơn?
	 \choice
		{$3$} 
		{$12$} 
		{\True  $60$}
		{$75$} 
	\loigiai{ 
		Có $5$ cách chọn $1$ món ăn trong $5$ món ăn, $4$ cách chọn $1$ loại quả tráng miệng trong $4$ loại quả  tráng miệng và $3$ cách chọn $1$ loại nước uống trong $3$ loại nước uống.\\
		Theo  quy tắc nhân có $5\cdot 4\cdot3=60$ cách chọn thực đơn.
	 }
\end{ex}

% nguồn: MaHoaNH_TracNghiem_Lop10
\begin{ex}%[0D8N1-2]%[ma:5d2e2fcd]
	 Có $3$ kiểu mặt đồng hồ đeo tay (vuông, tròn, elip) và $4$ kiểu dây (kim loại, da, vải và nhựa). Hỏi có bao nhiêu cách chọn một chiếc đồng hồ gồm một mặt và một dây?
	 \choice
		 {$7$}
		{$4$} 
		 {\True $12$}
		{$16$} 
	 \loigiai{
		Chọn $1$ kiểu mặt từ $3$ kiểu  mặt có $3$ cách.\\
		Chọn $1$ kiểu dây từ $4$  kiểu dây có $4$ cách.\\
		Vậy theo quy tắc  nhân có $12$ cách chọn $1$ chiếc đồng hồ gồm một mặt và một dây.
	} 
\end{ex}

% nguồn: MaHoaNH_TracNghiem_Lop10
\begin{ex}%[0D8N1-2]%[ma:9402fef9]
	Một khu di tích nọ có bốn cửa Đông, Tây, Nam, Bắc. Một người đi vào tham quan rồi đi ra, khi vào và ra phải đi hai cửa khác nhau. Tất cả các  cách đi vào và đi ra của người đó là
	 \choice
		 {$8$}
		{$4$} 
		 {\True $12$}
		 {$16$}
	\loigiai{ 
		Cứ mỗi cách đi vào sẽ có $3$  cách để đi ra.\\
		 Người đó có $4$ lựa chọn để đi vào.\\
		Do  đó có $4\cdot 3=12$ cách để đi vào và đi ra thỏa mãn yêu cầu bài toán.
	 }
\end{ex}

% nguồn: MaHoaNH_TracNghiem_Lop10
\begin{ex}%[0D8N1-2]%[ma:47a49f7e]
	 Một hộp đựng $5$ bi đỏ và $4$ bi xanh. Có bao nhiêu cách lấy $2$ bi có đủ cả $2$ màu?
	 \choice
		{\True $20$} 
		 {$9$}
		 {$36$}
		 {$16$}
	\loigiai{ 
		 Lấy $1$ bi đỏ có $5$ cách.\\
		Lấy  $1$ bi xanh có $4$ cách.\\
		Theo quy tắc nhân, số cách lấy $2$ bi  có đủ cả $2$ màu là $5\cdot 4=20$ cách.
	} 
\end{ex}

% nguồn: MaHoaNH_TracNghiem_Lop10
\begin{ex}%[0D8N1-2]%[ma:2ccbc8fc]
	Khối $10$ của một trường THPT có $280$ học sinh nam và $325$ học sinh nữ. Nhà trường cần chọn một học sinh ở khối $10$ đi dự dạ hội của học sinh thành  phố. Hỏi nhà trường có bao nhiêu cách chọn?
	\choice 
		{\True $605$} 
		 {$325$}
		{$280$} 
		{$91000$} 
	\loigiai{ 
		Số  cách chọn $1$ học sinh nam là $280$ cách.\\
		Số cách  chọn $1$ học sinh nữ là $325$ cách.\\
		Theo quy tắc cộng, số cách chọn $1$ học sinh bất kỳ  là $280+325=605$ (cách).
	} 
\end{ex}

% nguồn: MaHoaNH_TracNghiem_Lop10
\begin{ex}%[0D8N1-2]%[ma:c2c06c8e]
	Có bao nhiêu cách chọn $2$ số tự nhiên nhỏ hơn $7$, trong  đó có $1$ số lẻ và $1$ số chẵn?
	 \choice
		{\True  $12$}
		{$20$} 
		 {$9$}
		{$6$} 
	\loigiai{ 
		Tập các số tự nhiên nhỏ hơn $7$ là $\{0, 1, 2, 3, 4, 5, 6\}$. 
		\begin{itemize}
			\item Chọn $1$  số lẻ trong $3$ số lẻ có $3$ cách.
			\item Chọn $1$  số chẵn trong $4$ số chẵn có $4$ cách.
		\end{itemize}
		Áp dụng quy tắc nhân, có  $3\cdot 4=12$ cách.
	 }
\end{ex}

% nguồn: MaHoaNH_TracNghiem_Lop10
\begin{ex}%[0D8N1-2]%[ma:16402e7f]
	 Có bao nhiêu số tự nhiên có $5$ chữ số và các chữ số cách đều chữ số chính giữa là giống nhau?
	 \choice
		 {$9000$}
		 {\True $900$}
		 {$90000$}
		 {$500$}
	 \loigiai{
		Gọi số tự nhiên cần  tìm là $\overline{abcba}$ trong đó $a$, $b$, $c\in\mathbb{N},\,a\ne 0$.
		Ta có 
		\begin{itemize}
			\item $a$ có $9$ cách chọn  khác $0$.
			\item $b$ có $10$ cách  chọn.
			\item $c$ có  $10$ cách chọn.
		\end{itemize}
		Vậy có tất cả  là $9\cdot 10\cdot 10\cdot 1\cdot 1=900$ số.
	 }
\end{ex}

% nguồn: MaHoaNH_TracNghiem_Lop10
\begin{ex}%[0D8N1-2]%[ma:eb592aab]
	Một đội học sinh giỏi của trường THPT gồm 5 học sinh khối 12, 4 học sinh  khối 11, 3 học sinh khối 10.
	Số cách chọn ba học sinh trong đó mỗi khối có một  em là
	\choice 
		 {$12$}
		{$3$} 
		{$220$} 
		{\True  $60$}
	\loigiai{ 
		Để chọn một học sinh từ mỗi khối, ta  có
		\begin{itemize}
			\item Số cách chọn 1 học sinh từ khối 12 là  $C^1_5 = 5 $.
			\item Số cách chọn 1 học sinh từ khối 11 là $  C^1_4 = 4 $.
			\item Số cách chọn 1 học sinh từ khối 10 là  $ C^1_3 = 3 $.
		\end{itemize}
		Theo quy tắc nhân. Tổng số cách chọn là $C^1_5 \cdot  C^1_4 \cdot C^1_3 = 60.$
	 }
\end{ex}

% nguồn: MaHoaNH_TracNghiem_Lop10
\begin{ex}%[0D8N1-2]%[ma:75e1c653]
	Giả sử từ tỉnh A đến tỉnh B có thể đi bằng các phương tiện: ô tô, tàu hỏa, tàu thủy hoặc máy bay. Mỗi ngày có $10$ chuyến ô tô, $5$ chuyến tàu hỏa, $3$ chuyến tàu thủy và $2$ chuyến máy bay. Hỏi có  bao nhiêu cách đi từ tỉnh A đến tỉnh B?
	 \choice
		{$20$} 
		{\True  $300$}
		{$18$} 
		 {$15$}
	\loigiai{ 
		Vì từ tỉnh A đến tỉnh B có thể đi bằng các phương tiện: ô tô, tàu hỏa, tàu thủy hoặc  máy bay và
		\begin{itemize}
			 \item ôtô có $10$ cách chọn,
			\item tàu hỏa có  $5$ cách chọn,
			\item tàu thủy có $3$ cách chọn, 
			\item máy bay có $2$ cách  chọn.
		\end{itemize}
		Vậy theo quy tắc nhân, ta có $10\cdot 5\cdot 3\cdot 2 = 300$ cách  đi từ tỉnh A đến tỉnh B.
	 }
\end{ex}

% nguồn: MaHoaNH_TracNghiem_Lop10
\begin{ex}%[0D8H1-2]%[ma:a2f43ec6]
	Có bao nhiêu cách xếp $5$ bạn An,  Bình, Chi, Dũng, Lệ ngồi vào một chiếc ghế dài có $5$ chỗ ngồi sao cho An và Dũng luôn ngồi ở hai đầu ghế?
	 \choice
		{\True  $12$}
		 {$24$}
		 {$16$}
		{$120$} 
	\loigiai{ Xếp An,  Dũng ngồi hai đầu ghế có $2!=2$ cách xếp.\\
		Xếp Bình, Chi,  Lệ vào ba vị trí còn lại có $3!=6$ cách xếp.\\
		Vậy  có $2\cdot 6=12$ cách xếp thỏa yêu cầu bài toán.
	} 
\end{ex}

% nguồn: MaHoaNH_TracNghiem_Lop10
\begin{ex}%[0D8H1-2]%[ma:811df87a]
	Số điện thoại ở Huyện Củ  Chi có $7$ chữ số và bắt đầu bởi $3$ chữ số đầu tiên là $790$. Hỏi ở Huyện Củ Chi có tối đa bao nhiêu số điện thoại?
	\choice 
		{$1\,000$} 
		 {$100\,000$}
		 {\True $10\,000$}
		 {$1\,000\,000$}
	 \loigiai{
		Gọi số  điện thoại cần tìm có dạng $\overline{790abcd}$.\\
		Khi đó, ta có $10$ cách chọn  $a$, $10$ cách chọn $b$, $10$ cách chọn $c$, $10$ cách chọn $d$.\\
		Vậy số các số  điện thoại tối đa là $10^4$ số.
	} 
\end{ex}

% nguồn: MaHoaNH_TracNghiem_Lop10
\begin{ex}%[0D8H1-2]%[ma:14e5868b]
	Nhãn  mỗi chiếc ghế trong hội trường gồm hai phần: phần đầu là một chữ cái trong bảng chữ cái tiếng Anh, phần thứ hai là một số nguyên dương nhỏ hơn $26$. Hỏi có nhiều nhất bao nhiêu chiếc ghế được ghi nhãn khác nhau?
	\choice 
		 {$624$}
		{$48$} 
		 {\True $650$}
		 {$625$}
	\loigiai{ 
		Phần đầu là một chữ cái tiếng Anh nên  ta có $26$ cách chọn phần đầu.\\
		Phần thứ hai là một số nguyên  dương nhỏ hơn $26$ nên ta có $25$ cách chọn phần thứ hai.\\
		Vậy theo  qui tắc nhân ta có $26\cdot25=650$ cách.
	 }
\end{ex}

% nguồn: MaHoaNH_TracNghiem_Lop10
\begin{ex}%[0D8H1-2]%[ma:ce16ecea]
	Số các số lẻ có hai chữ số khác nhau  là
	 \choice
		 {$10$}
		 {$20$}
		{$30$} 
		{\True  $40$}
	\loigiai{ 
		Gọi số  cần tìm có dạng $\overline{a_1 a_2},\, \left(a_1\ne a_2\right)$.\\
		$a_1$ có $5$  cách chọn.\\
		$a_2$ có $8$ cách  chọn.\\
		Vậy  có $40$ số.
	 }
\end{ex}

% nguồn: MaHoaNH_TracNghiem_Lop10
\begin{ex}%[0D8H1-2]%[ma:10a024ef]
	 Số các số nguyên gồm $3$ chữ số khác nhau là
	\choice 
		 {$810$}
		{\True  $648$}
		 {$729$}
		 {$720$}
	 \loigiai{
		Gọi $n=\overline{a_1 a_2 a_3}$.\\ 
		 $a_1$ có $9$ cách chọn.\\
		$a_2$ có $9$ cách  chọn.\\
		 $a_3$ có $8$ cách chọn.\\
		Số cần tìm $9\times 9\times 8=648$. 
	} 
\end{ex}

% nguồn: MaHoaNH_TracNghiem_Lop10
\begin{ex}%[0D8H1-2]%[ma:ba588dd4]
	Từ các chữ số $1$,  $2$, $3$, $4$ có thể lập được bao nhiêu số có $3$ chữ số khác nhau?
	\choice 
		{$6$} 
		 {$8$}
		 {$12$}
		 {\True $24$}
	 \loigiai{
		Gọi số cần tìm có dạng $\overline{a_1  a_2 a_3}$.\\
		$a_1$ có $4$ cách chọn.\\ 
		$a_2$ có $3$ cách  chọn.\\
		$a_3$ có $2$ cách chọn.\\ 
		Vậy  có $24$ số.
	} 
\end{ex}

% nguồn: MaHoaNH_TracNghiem_Lop10
\begin{ex}%[0D8H1-2]%[ma:f5a194c1]
	Từ tập $X=\left\{1,2,3\right\}$ có thể lập được bao nhiêu số có $5$ chữ số mà chữ số $1$ có mặt đúng $3$ lần, còn các  chữ số khác có mặt đúng $1$ lần?
	\choice 
		 {$60$}
		{$10$} 
		 {\True $20$}
		{$30$} 
	 \loigiai{
		Xem số  phải lập như một dãy có $5$ ô trống.\\
		Đặt số $2$ vào có  $5$ cách.\\
		Đặt số $3$  vào có $4$ cách.\\
		Đặt $3$ chữ số $1$ vào có $1$  cách.\\
		 Vậy có $20$ cách.
	 }
\end{ex}

% nguồn: MaHoaNH_TracNghiem_Lop10
\begin{ex}%[0D8H1-2]%[ma:9c27dc52]
	Trong một buổi khiêu vũ có $ 20 $ nam và $ 18 $ nữ. Hỏi có bao nhiêu cách chọn ra một đôi  nam nữ để khiêu vũ?
	 \choice
		 {$703$}
		 {\True $360$}
		{$1406$} 
		 {$38$}
	 \loigiai{
		\begin{itemize}
			\item Chọn một nam có $ 20 $ cách. 
			\item Chọn một nữ  có $ 18 $ cách.
		\end{itemize}
		 Theo quy tắc nhân có $ 20\cdot 18=360 $ cách chọn ra một đôi nam nữ.
	 }
\end{ex}

% nguồn: MaHoaNH_TracNghiem_Lop10
\begin{ex}%[0D8H1-2]%[ma:c599d886]
	Trong kinh doanh nhà hàng, combo là một  hình thức gọi món theo thực đơn được kết hợp từ nhiều món ăn hoặc đồ uống. Nếu nhà hàng có $5$ món rau, $6$ món cá và $7$ món thịt thì có bao nhiêu cách tạo ra $1$ combo? Biết mỗi combo có đầy đủ $1$ món rau, $1$ món cá và $1$ món thịt.
	\choice 
		 {$18 $}
		 {$18!$}
		{$\mathrm{C}_{18}^3$} 
		{\True  $210$}
	\loigiai{Mỗi cách chọn một combo được thực  hiện qua $3$ công đoạn
		\begin{itemize}
			\item  \textbf{Công đoạn 1}: Chọn $1$ món rau từ $5$ món rau thì có $5$ cách.
			\item  \textbf{Công đoạn 2}: Chọn $1$ món cá từ $6$ món rau thì có $6$ cách.
			\item \textbf{Công đoạn 3}:  Chọn $1$ món thịt từ $7$ món thịt thì có $7$ cách.
		\end{itemize}
	Theo  quy tắc nhân, số cách chọn một combo là $5\cdot 6\cdot 7=210$.}
\end{ex}

% nguồn: MaHoaNH_TracNghiem_Lop10
\begin{ex}%[0D8H1-2]%[ma:436449a0]
	Có $10000$ vé số được đánh số từ $0000$ đến $9999$. Số các vé có $4$ chữ số khác  nhau là
	\choice 
		{$30240$} 
		{\True $5040$} 
		 {$10000$}
		 {$2520$}
	\loigiai{ 
		Gọi số  intrên vé $n=\overline{a_1 a_2 a_3 a_4}$.\\
		$a_1$ có $10$  cách chọn.\\
		$a_2$ có $9$ cách chọn.\\ 
		$a_3$  có $8$ cách chọn.\\
		$a_4$ có $7$ cách  chọn.\\
		Vậy số cần tìm  là $5\,040$.
	} 
\end{ex}

% nguồn: MaHoaNH_TracNghiem_Lop10
\begin{ex}%[0D8H1-2]%[ma:3b31f023]
	Cho tập $A = \left\lbrace 0;1;2;3;4;5\right\rbrace$. Có bao nhiêu số tự nhiên lẻ  gồm $2$ chữ số được lập từ tập $A$?
	 \choice
		 {$9$}
		{$12$} 
		{\True  $15$}
		 {$18$}
	 \loigiai{
		Gọi $\overline{ab}$ là số tự nhiên lẻ gồm $2$ chữ số được lập từ tập $A$. Khi  đó
		\begin{itemize}
			\item $a\in A \setminus  \left\lbrace 0\right\rbrace $ nên $a$ có $5$ cách chọn,
			\item $b \in \left\lbrace 1;3;5 \right\rbrace $ nên $b$ có $3$  cách chọn.
		\end{itemize}
		Vậy theo quy tắc nhân, ta có $5\cdot 3 = 15$ số tự nhiên lẻ gồm $2$ chữ số được lập  từ tập $A$.
	} 
\end{ex}

% nguồn: MaHoaNH_TracNghiem_Lop10
\begin{ex}%[0D8H1-2]%[ma:e032a41c]
	Có bao nhiêu số tự  nhiên có $2$ chữ số mà cả hai chữ số đều là lẻ?
	 \choice
		 {\True $25$}
		 {$10$}
		{$20$} 
		 {$50$}
	\loigiai{ 
		 Gọi số cần lập là $\overline{ab}$.
		\begin{itemize}
			\item Chọn $a$  lẻ có $5$ cách.
			\item Chọn $b$  lẻ có $5$ cách.
		\end{itemize}
		Theo quy  tắc nhân, có $5\cdot 5=25$ cách lập.
	} 
\end{ex}

% nguồn: MaHoaNH_TracNghiem_Lop10
\begin{ex}%[0D8H1-2]%[ma:dc995a4c]
	Có  bao nhiêu số tự nhiên có hai chữ số mà hai chữ số đều chẵn?
	\choice 
		 {$99$}
		 {$50$}
		{\True $20$} 
		 {$10$}
	\loigiai{ 
		Gọi số cần tìm có  dạng $\overline{ab}$ với $(a, b)\in A=\left\{0,2, 4,6, 8\right\}$ và $a\ne 0$.\\
		Trong  đó
		\begin{itemize}
			\item $a$ được chọn từ tập  $A\setminus\{0\}$ (có $4$ phần tử) nên có $4$ cách chọn.
			\item $b$ được chọn từ tập $A$ (có $5$ phần tử) nên có $5$  cách chọn.
		\end{itemize}
		Như vậy, ta có $4\cdot 5=20$ số cần  tìm.
	} 
\end{ex}

% nguồn: MaHoaNH_TracNghiem_Lop10
\begin{ex}%[0D8H1-2]%[ma:80fd6609]
	Biển số xe máy của tỉnh $A$ có $6$ kí tự, trong đó kí tự ở vị trí đầu tiên là một chữ cái trong bảng chữ cái tiếng Anh, kí tự ở vị trí thứ hai là một chữ số thuộc tập $\left\{1;2;\ldots;9\right\}$, mỗi kí tự ở bốn vị trí tiếp theo là một chữ số thuộc tập $\left\{0;1;2;\ldots;9\right\}$. Hỏi nếu chỉ dùng một mã số tỉnh thì tỉnh $A$ có thể làm được nhiều nhất bao nhiêu biển số xe máy khác nhau? 
	\choice 
		{\True $2\,340\,000$} 
		 {$234\,000$}
		 {$75$}
		{$2\,600\,000$} 
	 \loigiai{
		Giả sử biển số xe có dạng $a_1  a_2 a_3 a_4 a_5 a_6$.\\
		Bảng chữ cái tiếng Anh có $26$ chữ nên ta có  $26$ cách chọn $a_1$.\\
		Có $9$  cách chọn $a_2$.\\
		 Có $10$ cách chọn $a_3$.\\
		Có  $10$ cách chọn $a_4$.\\
		Có  $10$ cách chọn $a_5$.\\
		Có $10$ cách chọn  $a_6$.\\
		Vậy theo qui tắc nhân ta có $26\times  9\times 10\times 10\times 10\times 10=2\,340\,000$ biển số xe.
	 }
\end{ex}

% nguồn: MaHoaNH_TracNghiem_Lop10
\begin{ex}%[0D8H1-2]%[ma:d1cf7933]
	 Từ các chữ số $0$, $1$, $2$, $3$, $4$, $5$ có thể lập được bao nhiêu số lẻ gồm $4$ chữ số khác nhau?
	\choice 
		{$154$} 
		 {$145$}
		 {\True $144$}
		{$155$} 
	 \loigiai{
		Gọi số cần tìm có dạng $\overline{abcd}$ với $\left\{a, b, c, d\right\}\in A=\left\{0,1, 2,3, 4,5\right\}$. 
		\begin{itemize}
			 \item Vì $\overline{abcd}$ là số lẻ $\Rightarrow d\in \left\{1,3, 5\right\} \Rightarrow d$ có $3$ cách chọn.
			\item $a$ có $4$  cách chọn (khác $0$ và $d$).
			\item $b$ có  $4$ cách chọn.
			\item $c$ có $3$ cách  chọn.
		\end{itemize}
		Vậy có tất cả $3\cdot 4\cdot 4\cdot  3=144$ số cần tìm.
	} 
\end{ex}

% nguồn: MaHoaNH_TracNghiem_Lop10
\begin{ex}%[0D8H1-2]%[ma:b36315ae]
	Trong một tuần bạn $A$ dự định mỗi ngày đi thăm một người bạn trong $12$ người bạn của mình. Hỏi bạn $A$  có thể lập được bao nhiêu kế hoạch đi thăm bạn của mình?
	 \choice
		 {\True $3\,991\,680$}
		{$12!$} 
		{$35\,831\,808$} 
		 {$7!$}
	\loigiai{ 
		Một tuần có bảy ngày và mỗi ngày thăm  một bạn.\\
		Có  $12$ cách chọn bạn vào ngày thứ nhất.\\
		Có $11$ cách chọn bạn  vào ngày thứ hai.\\
		Có $10$ cách chọn bạn  vào ngày thứ ba.\\
		Có $9$  cách chọn bạn vào ngày thứ tư.\\
		Có $8$ cách chọn  bạn vào ngày thứ năm.\\
		Có $7$  cách chọn bạn vào ngày thứ sáu.\\
		Có $6$ cách chọn bạn vào ngày  thứ bảy.\\
		Vậy theo qui tắc nhân ta có $12\times 11\times 10\times 9\times 8\times 7\times 6=3\,991\,680$  cách.
	} 
\end{ex}

% nguồn: MaHoaNH_TracNghiem_Lop10
\begin{ex}%[0D8H1-2]%[ma:c3d7549e]
	Từ các chữ số $1$, $5$, $6$, $7$ có thể lập được bao nhiêu chữ số tự nhiên có  $4$ chữ số khác nhau?
	 \choice
		{$36$} 
		 {\True $24$}
		 {$20$}
		{$14$} 
	\loigiai{ 
		 Gọi số cần tìm có dạng $\overline{abcd}$ với $(a,b,c,d)\in A=\left\{1,5,6,7\right\}$.\\
		Vì  số cần tìm có $4$ chữ số khác nhau nên
		\begin{itemize}
			\item $a$ được chọn từ tập $A$  (có $4$ phần tử) nên có $4$ cách chọn.
			 \item $b$ được chọn từ tập $A\setminus\{a\}$ (có $3$ phần tử) nên có $3$ cách chọn.
			\item $c$ được chọn từ tập  $A\setminus\left\{a,b\right\}$ (có $2$ phần tử) nên có $2$ cách chọn.
			\item $d$ được chọn từ tập  $A\setminus\left\{a,b,c\right\}$ (có $1$ phần tử) nên có $1$ cách chọn.
		\end{itemize}
		Như vậy ta  có $4\cdot 3\cdot 2\cdot 1=24$ số cần tìm.
	 }
\end{ex}

% nguồn: MaHoaNH_TracNghiem_Lop10
\begin{ex}%[0D8H1-2]%[ma:0698edf0]
	\immini{Mã xác thực (OTP – One Time Password) do một ngân hàng gửi vào điện thoại của khách hàng cho mỗi  lần giao dịch là một dãy $ 6 $ kí tự từ các chữ số từ $ 0 $ đến $ 9 $. Có thể tạo ra bao nhiêu mã xác thực khác nhau như vậy?
	\choice 
		 {$151200$}
		{$136080$} 
		 {\True $1000000$}
		 {$900000$}
		}{ 
		\begin{tabular}{|l|}
			 \hline
			\texttt{Mã OTP xác  thực giao } \\
			\texttt{dịch là $ 712892  $, hiệu}\\
			\texttt{lực trong 1 phút } $\ldots$\\ 
			 \hline
	\end{tabular}}
	 \loigiai{
		Có $  10 $ cách chọn chữ số cho mỗi kí tự của mã xác nhận.\\
	Do đó  theo quy tắc nhân, số mã xác nhận có thể tạo ra là $10^6=1000000$.}
\end{ex}

% nguồn: MaHoaNH_TracNghiem_Lop10
\begin{ex}%[0D8V1-2]%[ma:6fd9fb47]
	Trong một giải thi  đấu bóng đá có $20$ đội tham gia với thể thức thi đấu vòng tròn. Cứ hai đội thì gặp nhau đúng một lần. Hỏi có tất cả bao nhiêu trận đấu xảy ra.
	 \choice
		{\True $190$} 
		{$182$} 
		{$280$} 
		 {$194$}
	\loigiai{ 
		Mỗi đội sẽ đấu với $19$  đội còn lại. Do đó, $20$ đội sẽ phải đấu $19\cdot20$ trận.\\
		Tuy nhiên, với cách đếm như thế, ta đã đếm trận đội  $A$ đấu với đội $B$ hai lần nên số trận đấu thực sự là $\dfrac{19\cdot20}{2}=190$ trận.
	 }
\end{ex}

% nguồn: MaHoaNH_TracNghiem_Lop10
\begin{ex}%[0D8V1-2]%[ma:f8eaff0e]
	Một phiếu điều tra về đề tự học của học sinh gồm $10$ câu  hỏi trắc nghiệm, mỗi câu có bốn lựa chọn để trả lời. Khi tiến hành điều tra, phiếu thu lại được coi là hợp lệ nếu người được hỏi trả lời đủ $10$ câu hỏi, mỗi câu chỉ chọn một phương án. Hỏi cần tối thiểu bao nhiêu phiếu hợp lệ để trong số đó luôn có ít nhất hai phiếu trả lời giống hệt nhau cả $10$ câu hỏi?
	\choice 
		{$2\,097\,152$} 
		{$10\,001$} 
		{\True $1\,048\,577$} 
		 {$1\,048\,576$}
	\loigiai{ 
		Mỗi câu hỏi có $4$ phương án  nên $10$ câu hỏi có $4^{10}=1\,048\,576$ phương án trả lời khác nhau.\\
		Theo nguyên lí Dirichlet, cho $n$ con thỏ và $n+1$ cái chuồng, khi đó sẽ có ít  nhất một chuồng chứa hai con thỏ.\\
		Ta coi số phương án trả lời là $n=1\,048\,576$, khi đó số phiếu hợp lệ tối thiểu cần phát tương tự như số chuồng  thỏ là $n+1=1\,048\,577$ phiếu.
	} 
\end{ex}


%%%%%%%%%%%%%%%% ĐÚNG/SAI %%%%%%%%%%%%%%%%

% nguồn: MaHoaNH_DungSai_Lop10
\begin{ex}%[0D8N1-2]%[ma:773f8445]
	Có $3$ thành phố $A$,  $B$, $C$. Có $5$ con đường nối từ $A$ đến $B$, có $4$ con đường nối từ $B$ đến $C$.
	\choiceTF 
		{\True Có  $20$ cách chọn một con đường để đi từ $A$ qua $B$ rồi đến $C$}
		{Có $9$ cách chọn một  con đường để đi từ $A$ qua $B$ rồi đến $C$}
		{Không có con đường  nào để đi từ $A$ qua $B$ rồi đến $C$}
		{\True Có $20$ cách chọn một  con đường để đi từ $C$ qua $B$ rồi đến $A$}
	 \loigiai{
		Theo quy tắc nhân, ta có $5\cdot 4=20$ cách chọn một con đường để đi từ $A$ qua $B$ rồi đến $C$. Ngược lại, ta  cũng có $4\cdot 5=20$ cách chọn một con đường để đi từ $C$ qua $B$ rồi đến $A$.
	 }
\end{ex}

% nguồn: MaHoaNH_DungSai_Lop10
\begin{ex}%[0D8H1-2]%[ma:ddebc133]
	Tại một nhà hàng chuyên phục vụ cơm trưa văn phòng, thực đơn gồm có $6$ món chính, $5$ món phụ và $4$ loại đồ uống. Tại đây thực khách có thể  chọn bữa trưa của mình gồm $1$ món chính, $1$ món phụ, $1$ loại đồ uống.
	\choiceTF 
		{\True  Có $120$ cách chọn một bữa trưa}
		{Có $15$ cách chọn một bữa  trưa}
		{Có $30$ cách chọn một bữa trưa} 
		{Có $20$  cách chọn một bữa trưa}
	\loigiai{ 
		Theo quy tắc nhân, ta có $6\cdot 5\cdot 4=120$ cách chọn  một bữa trưa.
	 }
\end{ex}

% nguồn: MaHoaNH_DungSai_Lop10
\begin{ex}%[0D8H1-2]%[ma:f2e3329c]
	Một đội học sinh giỏi của trường THPT gồm $5$ học  sinh khối $12$, $4$ học sinh khối $11$ và $3$ học sinh khối $10$.
	\choiceTF 
		{\True Có $60$ cách chọn  ba học sinh trong đó mỗi khối có một em}
		{Có $35$ cách chọn ba học sinh  trong đó có đúng một học sinh lớp $12$}
		{Có $30$ cách chọn  hai học sinh của đúng hai khối}
		{\True  Có $12$ cách chọn một học sinh giỏi của trường để phát biểu}
	\loigiai{ 
		\begin{itemchoice}
			 \itemch Có $5$ cách chọn học sinh khối $12$.\\
			Có $4$ cách chọn học sinh khối $11$.\\ 
			Có $3$ cách chọn học  sinh khối $10$.\\
			Vậy theo qui  tắc nhân ta có $5\times 4\times 3=60$ cách.
			\itemch  Có $5$ cách chọn học sinh khối $12$.\\
			Có $7$ học sinh của hai khối $10$ và  $11$, chọn học sinh thứ hai có $7$ cách chọn, chọn học sinh thứ ba có $6$ cách chọn.\\
			Với cách chọn lớp $10$  và $11$ như trên thì mỗi trường hợp lặp lại $2$ lần.\\
			 Vậy theo qui tắc nhân ta có $\dfrac{1}{2}\left( 5\times 7\times 6 \right)=105$ cách chọn.
			\itemch Nếu chọn 1 học sinh khối $12$ và $1$ học sinh khối $11$ có $5\cdot4=20$ cách  chọn.\\
			 Nếu chọn $1$ học sinh khối $11$ và $1$ học sinh khối $10$ có $4\cdot3=12$ cách chọn.\\
			Nếu chọn $1$ học  sinh khối $10$ và $1$ học sinh khối $12$ có $3\cdot5=15$ cách chọn.\\
			Theo qui  tắc cộng, ta có $20+12+15=47$ cách chọn.
			\itemch  Nếu chọn $1$ học sinh khối $12$ có $5$ cách chọn.\\
			 Nếu chọn $1$ học sinh khối $11$ có $4$ cách.
			Nếu chọn $1$ học sinh khối $10$ có  $3$ cách.
			Theo qui tắc cộng, ta có $5+4+3=12$  cách chọn.
		\end{itemchoice}
	} 
\end{ex}


%%%%%%%%%%%%%%%% TRẢ LỜI NGẮN %%%%%%%%%%%%%%%%

% nguồn: MaHoaNH_TraLoiNgan_Lop10
\begin{bt}%[0D8H1-2]%[ma:cffad701]
	Có $4$ bông hoa hồng khác nhau,  có $6$ bông hoa lan khác nhau, có $5$ bông hoa cúc khác nhau. Hỏi bạn có bao nhiêu cách chọn $3$ bông hoa để cắm sao cho hoa trong lọ phải có một bông hoa của mỗi loại?
	\shortans{120} 
	 \loigiai{
		Có $4$ cách chọn một bông hoa hồng, $6$ cách chọn  một bông hoa lan, $5$ cách chọn một bông hoa cúc để cắm vào lọ.\\
		 Theo quy tắc nhân, ta có số cách chọn $3$ bông hoa để cắm sao cho hoa trong lọ phải có một bông hoa của mỗi loại là $4\cdot6 \cdot5=120$ cách.
	} 
\end{bt}

% nguồn: MaHoaNH_TraLoiNgan_Lop10
\begin{bt}%[0D8H1-2]%[ma:1f39e8f3]
	Một thùng có $12$ hộp đựng bút màu đỏ, $18$ hộp đựng bút màu xanh. Số cách khác nhau để chọn được đồng thời một hộp màu đỏ, một hộp màu xanh  là bao nhiêu?
	 \shortans{216}
	 \loigiai{
		Có $12$ cách chọn một hộp  màu đỏ, có $18$ cách chọn một hộp màu xanh.\\
		Theo quy tắc nhân thì số cách khác nhau để chọn được đồng thời một hộp màu đỏ, một hộp màu  xanh là $12\cdot18=216$.
	 }
\end{bt}

% nguồn: MaHoaNH_TraLoiNgan_Lop10
\begin{bt}%[0D8H1-2]%[ma:abc849c5]
	Có bao nhiêu cách xếp $4$ người $A$, $B$, $C$, $D$ lên $3$ hàng ghế trên một toa tàu được đánh thứ tự $1$, $2$, $3$. Biết mỗi hàng  ghế có thể chứa tối đa $4$ người?
	\shortans{81} 
	 \loigiai{
		Xếp $A$ lên  một trong $3$ hàng ghế có $3$ cách.\\
		Xếp  $B$ lên một trong $3$ hàng ghế có $3$ cách.\\
		Tương tự, số cách xếp $C$ và $D$ cũng là $3$  cách.\\
		Với  mỗi cách xếp $A$ ta có $3$ cách xếp $B$ lên hàng ghế.\\
	Vậy số cách xếp thỏa mãn là $3\cdot 3\cdot 3\cdot 3=81$ cách.} 
\end{bt}

% nguồn: MaHoaNH_TraLoiNgan_Lop10
\begin{bt}%[0D8H1-2]%[ma:0226f064]
	Có $9$ cặp vợ chồng đi dự tiệc. Hỏi có bao nhiêu cách chọn một người đàn ông và một người đàn bà trong bữa tiệc phát  biểu ý kiến sao cho hai người đó không là vợ chồng?
	 \shortans{72}
	\loigiai{ 
		Để  chọn một người đàn ông và một người đàn bà trong bữa tiệc phát biểu ý kiến sao cho hai người đó không là vợ chồng ta tiến hành hai công đoạn liên tiếp.
		\begin{enumerate} [$\bullet$]
			\itemch \textbf{Công đoạn 1:}  Chọn $1$ người đàn ông trong $9$ người đàn ông lên phát biểu có $9$ cách chọn.
			\itemch \textbf{Công đoạn 2:} Chọn $1$ người đàn bà trong $9$ người đàn bà còn lại sao cho người đó không phải là vợ của người đàn ông đã  chọn thì có $8$ cách.
		\end{enumerate}
		Theo quy tắc nhân  thì có $9\cdot8=72$ cách.
	 }
\end{bt}

% nguồn: MaHoaNH_TraLoiNgan_Lop10
\begin{bt}%[0D8H1-2]%[ma:780309f9]
	Đội tuyển học sinh giỏi Toán gồm $10$ em: $5$ nam và $5$ nữ. Muốn chọn ra $1$  bạn nam làm tổ trưởng, $1$ bạn nữ làm tổ phó và $1$ thư ký. Số cách chọn là bao nhiêu?
	 \shortans{200}
	 \loigiai{
		\begin{itemize}
			\item Chọn $1$ bạn nam làm tổ trưởng có $5$ cách  chọn.
			\item Mỗi cách chọn $1$ bạn nam làm tổ trưởng có $5$ cách chọn $1$ bạn nữ làm tổ phó. 
			\item Mỗi cách chọn $1$  bạn nam làm tổ trưởng và $1$ bạn nữ làm tổ phó thì có $8$ cách chọn $1$ bạn làm thư ký.
		\end{itemize}
		 Theo quy tắc nhân ta có $5 \cdot 5 \cdot 8=200$ cách chọn.
	} 
\end{bt}

% nguồn: MaHoaNH_TraLoiNgan_Lop10
\begin{bt}%[0D8H1-2]%[ma:c0010892]
	Một người vào cửa hàng ăn, người đó chọn thực đơn gồm $1$ món ăn trong $5$ món, $1$ loại quả  tráng miệng trong $5$ loại quả tráng miệng và một loại nước uống trong $3$ loại nước uống. Có bao nhiêu cách chọn thực đơn?
	\shortans{75} 
	 \loigiai{
		\begin{itemize}
			\item Chọn $1$  món ăn trong $5$ món: Có $5$ cách chọn.
			\item Chọn $1$ loại quả tráng miệng  trong $5$ loại quả tráng miệng: Có $5$ cách chọn.
			\item Chọn $1$ loại nước uống  trong $3$ loại nước uống: Có $3$ cách chọn.
		\end{itemize}
		Theo quy tắc nhân, có $5\cdot5 \cdot 3=75$ cách chọn thực  đơn gồm $1$ món ăn, $1$ loại quả tráng miệng và $1$ loại nước uống.
	} 
\end{bt}

% nguồn: MaHoaNH_TraLoiNgan_Lop10
\begin{bt}%[0D8H1-2]%[ma:3d984dc3]
	Một nhóm gồm $5$ em học sinh  (trong đó có một bạn tên Tùng) đang đứng xếp thành một hàng dọc, hỏi có bao nhiêu cách xếp sao cho Bạn Tùng đứng đầu hàng?
	\shortans{24} 
	 \loigiai{
		Giai đoạn 1: Xếp bạn Tùng đứng  ở đầu hàng có $1$ cách.\\
		Giai  đoạn 2: Chọn một học sinh đứng vị trí tiếp theo có $4$ cách.\\
		Giai đoạn 3: Chọn một học sinh đứng vị trí tiếp theo có  $3$ cách.\\
		Giai đoạn 4: Chọn một học sinh đứng vị trí tiếp theo có $2$ cách.\\ 
		Giai  đoạn 5: Chọn một học sinh đứng cuối hàng có $1$ cách.\\
		Số cách xếp một hàng thỏa mãn là  $1\cdot 4\cdot 3\cdot 2\cdot 1=24$ cách.
	} 
\end{bt}

% nguồn: MaHoaNH_TraLoiNgan_Lop10
\begin{bt}%[0D8H1-2]%[ma:ffe8d9b8]
	Một bó có $8$ hoa hồng trắng, $7$ hoa hồng đỏ và $10$ hoa hồng vàng. Hỏi có mấy cách chọn lấy ba bông  hoa có đủ cả ba màu?
	\shortans{560} 
	\loigiai{ 
		Để chọn ba bông hoa có đủ ba màu (nghĩa là $1$ hoa hồng trắng, $1$ hoa hồng đỏ và $1$ hoa  hồng vàng), ta có
		\begin{enumerate} [$\bullet$]
			\itemch Có $8$ cách chọn hoa  hồng trắng.
			\itemch  Có $7$ cách chọn hoa hồng đỏ.
			\itemch  Có $10$ cách chọn hoa hồng vàng.
		\end{enumerate}
		Vậy theo qui tắc nhân  ta có $8\times 7\times 10=560$ cách.
	 }
\end{bt}

% nguồn: MaHoaNH_TraLoiNgan_Lop10
\begin{bt}%[0D8H1-2]%[ma:05c7dc33]
	Từ các số $0$, $1$, $2$, $3$, $5$ có thể lập được bao nhiêu số tự nhiên không chia hết cho $5$ gồm  $4$ chữ số khác nhau?
	\shortans{54} 
	 \loigiai{
		Gọi số cần tìm có dạng  $\overline{a b c d}$.
		\begin{itemize}
			\item $d$ có $3$ cách chọn,  $d \notin\{0,5\}$.
			\item $a$ có $3$ cách chọn, $a  \neq 0$, $a \neq d$.
			\item $b$ có $3$ cách chọn, $b \neq a$, $b  \neq d$.
			\item $c$ có  $2$ cách chọn, $c \neq a$, $c \neq d$, $c \neq b$.
		\end{itemize}
		Suy ra  có $3\cdot3\cdot3\cdot2 =54$ số các số tự nhiên có $4$ chữ số khác nhau được lập từ các số $0$, $1$, $2$, $3$, $5$ không chia hết cho $5$ .
	 }
\end{bt}

% nguồn: MaHoaNH_TraLoiNgan_Lop10
\begin{bt}%[0D8H1-2]%[ma:97338c2a]
	Cần xếp $3$ nam, $3$ nữ vào $1$ hàng có $6$ ghế  (mỗi bạn một ghế). Hỏi có bao nhiêu cách xếp sao cho nam nữ ngồi xen kẽ?
	\shortans{72} 
	\loigiai{ 
		\begin{itemize}
			\item Có  $6$ cách chọn một người tuỳ ý ngồi vào chỗ thứ nhất.
			\item Tiếp đến, có $3$ cách chọn một  người khác phái ngồi vào chỗ thứ $2$.
			\item  Lại có $2$ cách chọn một người khác phái ngồi vào chỗ thứ $3$.
			\item Có $2$ cách chọn  vào chỗ thứ $4$.
			\item Có $1$  cách chọn vào chỗ thứ $5$.
			\item Có $1$ cách chọn vào chỗ thứ $6$. 
		\end{itemize}
		 Vậy có $6\cdot3\cdot2\cdot2\cdot1\cdot1=72$ cách.
	} 
\end{bt}

% nguồn: MaHoaNH_TraLoiNgan_Lop10
\begin{bt}%[0D8H1-2]%[ma:18f13368]
	Cho $n$ là  số tự nhiên. Xét các mệnh đề:\\
	$P$: \lq\lq $n$ là một số tự nhiên chia hết cho  $16$\rq\rq.\\
	$Q$:  \lq\lq $n$ là một số tự nhiên chia hết cho $8$\rq\rq.\\
	Cho biết  có bao nhiêu mệnh đề đúng trong các mệnh đề \lq\lq $P \Rightarrow Q$\rq\rq, \lq\lq $Q \Rightarrow P$\rq\rq\ và \lq\lq $P \Leftrightarrow Q$\rq\rq?
	\shortans{1} 
	\loigiai{ 
		\begin{itemize}
			\item Phát biểu mệnh đề $P \Rightarrow Q$: \lq\lq $n$ là một số tự  nhiên chia hết cho $16$ thì $n$ là một số tự nhiên chia hết cho $8$\rq\rq\\
			Mệnh đề này đúng, vì $n$ chia hết cho $16$ thì $n = 16k = 8 \cdot (2k)$, với $k \in  \mathbb{N}$ nên $n$ chia hết cho $8$.
			\item Phát  biểu mệnh đề $Q \Rightarrow P$: \lq\lq $n$ là một số tự nhiên chia hết cho $8$ thì $n$ là một số tự nhiên chia hết cho $16$\rq\rq.\\
			Mệnh đề  này sai, với $n = 8$ chia hết cho $8$ nhưng không chia hết cho $16$.
			\item Mệnh đề $P \Rightarrow Q$ đúng và  mệnh đề $Q \Rightarrow P$ sai nên $P \Leftrightarrow Q$ sai.
		\end{itemize}
		Vậy chỉ có một mệnh đề đúng. 
	} 
\end{bt}

% nguồn: MaHoaNH_TraLoiNgan_Lop10
\begin{bt}%[0D8H1-2]%[ma:c1839475]
	Cho  các mệnh đề\\
	$A$: \lq\lq Nếu  $\triangle ABC$ đều có cạnh bằng $a$, đường cao là $h$ thì $h = \dfrac{a\sqrt{3}}{2}$\rq\rq.\\
	$B$: \lq\lq Tứ giác có bốn cạnh bằng  nhau là hình vuông\rq\rq.\\
	$C$:  \lq\lq $15$ là số nguyên tố\rq\rq.\\
	$D$: \lq\lq $\sqrt{125}$ là một số  nguyên\rq\rq.\\
	Hãy cho biết trong các mệnh đề kéo theo sau có bao  nhiêu mệnh đề sai?\\
	$A \Rightarrow B$, $B \Rightarrow C$, $A \Rightarrow  D$.
	\shortans{2} 
	\loigiai{ 
		\begin{itemize}
			\item Mệnh đề $A \Rightarrow B$ là sai do  $A$ đúng nhưng $B$ sai.
			\item Mệnh đề $B \Rightarrow  C$ là đúng vì cả $B$ và $C$ đều sai.
			\item Mệnh đề  $A \Rightarrow D$ là sai do $A$ đúng nhưng $D$ sai.
		\end{itemize}
		 Vậy có hai mệnh đề sai.
	} 
\end{bt}

% nguồn: MaHoaNH_TraLoiNgan_Lop10
\begin{bt}%[0D8H1-2]%[ma:0b1e8d04]
	Ếch hay cóc?\\ 
	Trong một đầm lầy ma thuật, có  hai loài lưỡng cư biết nói: cóc luôn luôn nói đúng và ếch luôn luôn nói sai.\\
	Bốn loài lưỡng cư, Brian, Chris, LeRoy và Mike sống cùng nhau trong đầm lầy này và chúng đưa ra những tuyên  bố sau
	\begin{itemize}
		\item Brian: \lq \lq Mike và tôi là những loài khác nhau\rq  \rq.
		\item Chris:  \lq \lq LeRoy là một con ếch\rq \rq.
		\item LeRoy: \lq \lq Chris là một con ếch\rq  \rq.
		\item Mike: \lq \lq Trong bốn người  chúng tôi, ít nhất hai người là cóc\rq \rq.
	\end{itemize}
	 Có bao nhiêu loài lưỡng cư là ếch?
	\shortans  {$3$}
	 \loigiai{
		Giả sử Brian là cóc (nói thật) $\Rightarrow$ Mike là  ếch (nói dối).\\
		$\Rightarrow$ Chỉ có $1$ con là ếch  trong $4$ con mà Mike đã là ếch.\\
		$\Rightarrow$ LeRoy và Chris là đều cóc (nói  thật).\\
		Nhưng Chris nói LeRoy là ếch $\Rightarrow$ mâu  thuẫn.\\
		Vậy Brian nói dối (là Ếch) $\Rightarrow$ Brian và Mike cùng là loài ếch (nói dối).\\ 
		$\Rightarrow$ Chỉ có  1 con cóc và 3 con còn lại là ếch (*).\\
		Nếu Chris là Cóc (nói thật) $\Rightarrow$ LeRoy  là ếch (nói dối) $\Rightarrow$ Thỏa mãn (*).\\
		Nếu LeRoy là Cóc  (nói thật) $\Rightarrow$ Chris là ếch (nói dối) $\Rightarrow$ Thỏa mãn (*).\\
		Vậy có $3$ loài lưỡng cư  là ếch.
	 }
\end{bt}

% nguồn: MaHoaNH_TraLoiNgan_Lop10
\begin{bt}%[0D8V1-2]%[ma:90b707cd]
	Từ các  chữ số $1$, $2$, $3$, $4$, $5$, $6$, $7$, $8$ lập được bao nhiêu số tự nhiên có $4$ chữ số và chia hết cho $15$?
	 \shortans{171}
	 \loigiai{
		Trong $8$ chữ số đã cho có $2$ chữ  số chia hết cho $3$, $3$ chữ số chia $3$ dư $1$ và $3$ chữ số chia $3$ dư $2$ là $A=\{3;6\}$, $B=\{1;4;7\}$ và $C=\{2;5;8\}$.\\
		Gọi số  thỏa mãn yêu cầu bài toán có dạng $X=\overline{a b c d}\,\,(1 \leq a, b, c, d \leq 9)$.\\
		Để $X\vdots 15 \Rightarrow X\vdots  3$ và $X\vdots 5$.\\
		 $X\vdots 5 \Rightarrow d=5$.\\
		$X\vdots 3 \Rightarrow (a+b+c+5) \vdots 3  \Rightarrow a+b+c$ chia $3$ dư $1$.\\
		 Ta xét $3$ trường hợp:
		\begin{enumerate}[1.]
			\item Trường hợp $1$: $a+b$ chia hết cho $3$ và $c$ chia  $3$ dư $1$ thì có $3$ cách chọn $c$.\\
			Nếu $a$  và $b$ cùng chia hết cho $3$ có $2\cdot 2=4$ cách.\\
			Nếu $a$ chia $3$ dư $1$, $b$ chia cho $3$ dư $2$ có: $3\cdot 3 \cdot 2=18$ cách.\\ 
			Suy ra có $(4+18) \cdot  3=66$ số.
			\item Trường hợp $2$: $a+b$ chia $3$ dư $1$ và $c$  chia hết cho $3$ thì có $2$ cách chọn $c$.\\
			Nếu $a$ và $b$ cùng chia $3$ dư $2$ có: $3\cdot 3=9$ cách.\\ 
			Nếu $a$ chia $3$ dư $1$, $b$ chia hết cho $3$  có $3\cdot 2 \cdot 2=12$ cách.\\
			Suy ra có $(9+12)  \cdot 2=42$ số.
			\item Trường hợp $3$: $a+b$ chia $3$ dư $2$ và $c$ chia $3$  dư $2$ thì có $3$ cách chọn $c$.\\
			Nếu $a$ và $b$  cùng chia $3$ dư $1$ có $3\cdot 3=9$ cách.\\
			Nếu $a$ chia $3$ dư $2$, $b$ chia hết cho $3$ có: $3\cdot 2 \cdot  2=12$ cách.\\
			Suy ra có $(9+12) \cdot 3=63$ số. 
		\end{enumerate}
		Vậy  có $66+42+63=171$ số.
	} 
\end{bt}


%%%%%%%%%%%%%%%% TỰ LUẬN %%%%%%%%%%%%%%%%

% nguồn: MaHoaNH_TuLuan_Lop10
\begin{bt}%[0D8H1-2]%[ma:ebf0aa08]
	Có $3$ học sinh nữ và $4$ học sinh nam cùng xếp theo một hàng ngang. Có bao nhiêu cách xếp hàng để học sinh  nam và nữ xếp xen kẽ nhau?
	 \loigiai{
		 Hàng được xếp phải thỏa mãn Nam-Nữ-Nam-Nữ-Nam-Nữ-Nam.
		\begin{itemize}
			\item Chọn một  nam sinh cho vị trí thứ nhất có $4$ cách.
			\item Chọn một nữ sinh cho vị trí  thứ hai có $3$ cách.
			\item Số cách chọn học sinh cho các vị trí tiếp theo lần lượt là $3$, $2$, $2$,  $1$.
		\end{itemize}
		Vậy số cách xếp thỏa mãn là  $4\cdot 3\cdot 3\cdot 2\cdot 2\cdot 1=144$ cách.
	} 
\end{bt}

% nguồn: MaHoaNH_TuLuan_Lop10
\begin{bt}%[0D8H1-2]%[ma:3fb62717]
	Một hộp có $12$ viên bi xanh được đánh số từ $1$ đến $12$, $10$ viên bi đỏ được đánh số từ $1$ đến $10$ và $8$ viên bi vàng được đánh số từ $1$ đến $8$. Hỏi có bao nhiêu cách chọn ra $5$ viên bi bao gồm $1$ viên bi vàng, $2$ viên bi  đỏ và $2$ viên bi xanh sao cho $5$ viên bi được chọn mang $5$ số đôi một khác nhau?
	\loigiai{ 
		\begin{itemize}
			\item Số cách  chọn $1$ viên bi vàng bằng $\mathrm{C}_8^1$.
			\item Số cách chọn $2$ viên bi đỏ bằng $\mathrm{C}_{9}^2$. 
			\item Số  cách chọn $2$ viên bi xanh bằng $\mathrm{C}_{9}^2$.
		\end{itemize}
		Theo quy tắc nhân, số cách chọn ra $5$ bi thỏa yêu cầu bài toán  là
		 \[
		 \mathrm{C}_8^1\cdot \mathrm{C}_{9}^2\cdot \mathrm{C}_{9}^2=10368
		 \]
	} 
\end{bt}

% nguồn: MaHoaNH_TuLuan_Lop10
\begin{bt}%[0D8H1-2]%[ma:072803cc]
	Ba anh em An, Bình, Vinh ngồi làm  bài xung quanh một cái bàn được trải khăn mới. Khi phát hiện có vết mực, bà hỏi thì các cháu lần lượt trả lời
	\begin{itemize}
		\item An: \lq \lq Em Vinh không làm đổ mực,  đấy là do em Bình\rq \rq.
		\item Bình: \lq \lq  Em Vinh làm đổ mực, anh An không làm đổ mực\rq \rq.
		\item Vinh: \lq \lq Theo cháu, Bình không làm đổ mực,  còn cháu hôm nay không chuẩn bị bài\rq \rq.
	\end{itemize}
	Biết rằng trong $3$ em thì có $2$ em nói đúng, $1$ em nói  sai. Hỏi ai làm đổ mực?
	\loigiai{ 
		Nếu An nói đúng thì  Bình là người làm đổ, suy ra Bình nói sai, theo đề bài ta có Vinh nói đúng.\\
		Nếu Vinh nói đúng thì Bình không làm đổ mực.  Suy ra mâu thuẫn.\\
		Nếu Bình nói đúng, Vinh  làm đổ mực thì An nói sai.\\
		Dẫn đến Vinh nói đúng.  Suy ra thỏa mãn.\\
		Vậy Vinh  làm đổ mực.
	 }
\end{bt}

% nguồn: MaHoaNH_TuLuan_Lop10
\begin{bt}%[0D8H1-2]%[ma:0e17b227]
	Trên một hòn đảo, tôi đã gặp ba người $A$, $B$ và $C$, một người là hiệp sĩ, một người khác là kẻ bất lương và người kia là gián điệp. Người hiệp sĩ luôn nói sự thật, kẻ bất lương luôn nói dối và gián điệp có thể nói dối hoặc  nói thật.
	\begin{itemize}
		\item $A$  nói: \lq\lq Tôi là hiệp sĩ\rq\rq.
		\item $B$ nói: \lq\lq Tôi là kẻ  bất lương\rq\rq.
		\item $C$ nói: \lq\lq Tôi là gián  điệp\rq\rq.
	\end{itemize}
	Hỏi ai  là gián điệp?
	 \loigiai{
		 Ta xét các khả năng:
		\begin{itemize}
			\item Nếu $A$ nói thật  $\Rightarrow$ $A$ là hiệp sĩ $\Rightarrow$ $B$ hoặc $C$ là kẻ bất lương.
			\item Nếu $B$ là  kẻ bất lương thì đang nói thật $\Rightarrow$ Mâu thuẫn.
			\item Nếu $C$ là kẻ bất lương  $\Rightarrow$ $C$ nói dối $\Rightarrow$ không mâu thuẫn.
		\end{itemize}
		Như vậy, $A$ là hiệp sĩ (nói thật), $C$ là kẻ bất lương (nói dối),  $B$ là gián điệp.\\
		Vậy người gián điệp  là $B$.
	} 
\end{bt}

% nguồn: MaHoaNH_TuLuan_Lop10
\begin{bt}%[0D8V1-2]%[ma:3edb7b50]
	Bạn An có hộp chì màu gồm $12$ màu khác nhau. Bạn An muốn dùng $12$ viết chì màu đó để tô các cạnh của hình vuông $ABCD$  sao cho mỗi cạnh được tô $1$ màu và các cạnh có cùng đỉnh sẽ được tô $2$ màu khác nhau. Hỏi bạn An có bao nhiêu cách tô màu cạnh hình vuông $ABCD$ ?
	 \loigiai{
		\begin{enumerate}[1.]
			\item \textbf{Trường hợp 1:} Màu cạnh $CD$  trùng màu cạnh $AB$.
			\begin{enumerate}
				\item Số cách chọn màu cạnh $CD$: $11$ cách (trùng  màu cạnh $AB$)
				\item Số cách  chọn màu cạnh $DA$: $11$
				\item Tổng số cách tô màu:  $12 \cdot 11 \cdot 1 \cdot 11 = 15972$
			\end{enumerate}
			\item \textbf{Trường hợp 2:}  Màu cạnh $CD$ khác màu cạnh $AB$.
			\begin{enumerate}
				\item Số cách chọn  màu cạnh $CD$: $10$ (khác màu cạnh $AB$)
				\item Số cách chọn  màu cạnh $DA$: $10$
				\item Tổng số cách tô màu:  $12 \cdot 11 \cdot 10 \cdot 10 = 13200$
			\end{enumerate}
		\end{enumerate}
		Tổng số cách tô màu cạnh hình vuông $ABCD$: $15972 + 13200 =  29172$
	} 
\end{bt}

% nguồn: MaHoaNH_TuLuan_Lop10
\begin{bt}%[0D8Y1-2]%[ma:ac2ffb38]
	Ba lớp của một trường đang lên kế hoạch để đi  dã ngoại, mỗi lớp có thể chọn một trong năm địa điểm. Có bao nhiêu kết quả có thể xảy ra về cách chọn địa điểm của ba lớp?
	\loigiai{ 
	Mỗi lớp có 5 cách chọn địa điểm. Theo quy tắc nhân, số cách chọn địa điểm của ba lớp  là $5 \cdot 5 \cdot 5=125$.}
\end{bt}

% nguồn: MaHoaNH_TuLuan_Lop10
\begin{bt}%[0D8Y1-2]%[ma:faa658fb]
	Tung một đồng xu 5 lần liên tiếp và ghi lại kết quả (dùng  kí hiệu SSNSN để chỉ kết quả 5 lần tung lần lượt là sấp, sấp, ngửa, sấp, ngửa). Có bao nhiêu kết quả khác nhau có thể xảy ra?
	 \loigiai{
	Có thể coi việc tung đồng xu 5 lần liên tiếp là công việc gồm 5 công đoạn. Mỗi công đoạn có 2 phương án thực hiện, tương ứng đồng xu xuất hiện sấp hay ngửa. Do đó, theo quy tắc nhân, có $2 \cdot 2 \cdot 2 \cdot  2 \cdot 2=32$ kết quả có thể của việc tung đồng xu 5 lần liên tiếp.}
\end{bt}

% nguồn: MaHoaNH_TuLuan_Lop10
\begin{bt}%[0D8B1-2]%[ma:043749a5]
	Mã xác thực (OTP – One Time Password) do một ngân hàng gửi vào điện thoại của khách hàng cho mỗi lần giao dịch là một dãy 6 kí tự từ các chữ số từ 0 đến 9. Có thể tạo ra bao nhiêu mã  xác thực khác nhau như vậy?
	 \loigiai{
	Có 10 cách chọn chữ số cho mỗi kí tự của mã xác nhận. Do  đó theo quy tắc nhân, số mã xác nhận có thể tạo ra là $10^6=1000000$.}
\end{bt}

% nguồn: MaHoaNH_TuLuan_Lop10
\begin{bt}%[0D8Y1-2]%[ma:b86242a5]
	Tại một nhà hàng chuyên phục vụ cơm trưa văn phòng, thực đơn có $6$ món chính, $3$ món phụ và $3$ loại đồ uống. Thực khách có bao nhiêu cách lựa chọn bữa trưa gồm một món chính, một món phụ  và một loại đồ uống?
	\loigiai{ 
		Số cách lựa chọn bữa trưa gồm một món chính, một món phụ và một loại đồ uống là $6\cdot 3\cdot 3 = 54$  (cách).
	} 
\end{bt}

% nguồn: MaHoaNH_TuLuan_Lop10
\begin{bt}%[0D8B1-2]%[ma:2c5fd925]
	Dùng sáu chữ số 0; 1; 2; 3; 4;  5 có thể lập được bao nhiêu
	\begin{enumerate}
		\item mật khẩu có bốn chữ số khác  nhau?
		 \item số tự nhiên có bốn chữ số khác nhau?
		\item số tự nhiên lẻ có bốn chữ  số khác nhau?
	\end{enumerate}
	 \loigiai{
		\begin{enumerate}
			\item Kí hiệu mật khẩu cần lập là $abcd$, trong đó $a, b, c, d$ là các chữ số khác nhau từ sáu chữ số đã cho. Coi việc chọn mật khẩu là một công việc gồm bốn công đoạn.\\ 
			Công đoạn 1: Chọn chữ số $a$ từ sáu chữ số đã cho, có 6 cách  chọn.\\
			Công đoạn 2: Chọn chữ số $b$ từ năm chữ số còn lại,  có 5 cách chọn.\\
			 Công đoạn 3: Chọn chữ số $c$ từ bốn chữ số còn lại, có 4 cách chọn.\\
			Công đoạn  4: Chọn chữ số $d$ từ ba chữ số còn lại, có 3 cách chọn.\\
			Từ đó, theo quy tắc nhân, có thể lập được $6 \cdot 5 \cdot 4 \cdot 3=360$ mật  khẩu theo yêu cầu.
			\item Kí hiệu số tự nhiên cần lập là $\overline{abcd}$, trong đó $a, b, c, d$ là các chữ số khác nhau từ sáu chữ số đã cho, $a \ne 0$. Đầu tiên, có 5 cách chọn chữ số $a$. Tiếp theo, có 5 cách chọn chữ số $b$ từ các chữ số còn lại. Tiếp tục, có 4 cách chọn chữ số $c$ và 3 cách chọn chữ số $d$. Từ đó, theo quy tắc nhân, có $5 \cdot 5 \cdot 4 \cdot 3=300$ số  tự nhiên có bốn chữ số khác nhau lập từ các chữ số đã cho.
			\item Kí hiệu số tự nhiên lẻ có bốn chữ số khác nhau như câu b. Đầu tiên, có 3 cách chọn chữ số $d$ từ ba chữ số lẻ $1;3;5$. Tiếp theo, có 4 cách chọn chữ số $a$ từ 4 chữ số khác 0 còn lại. Tiếp theo, có 4 cách chọn chữ số $c$ từ 4 số còn lại. Cuối cùng, có 3 cách chọn chữ số $b$ từ các chữ  số còn lại. Từ đó, theo quy tắc nhân, có $3 \cdot 4 \cdot 4 \cdot 3=144$ số tự nhiên lẻ có bốn chữ số khác nhau lập từ các chữ số đã cho.
	\end{enumerate}}
\end{bt}

% nguồn: MaHoaNH_TuLuan_Lop10
\begin{bt}%[0D8T1-2]%[ma:5c1399ea]
	\immini{ 
	Một khoá tổ hợp với đĩa quay có 40 vạch số (xem Hình 7). Mật mã của khoá là một dãy gồm 3 số, kí hiệu là $a-b-c$, mỗi số là một số tự nhiên từ 0 đến 39. Để mở khoá, cần quay mặt  số ngược chiều kim đồng hồ cho đến khi điểm mốc gặp vạch số $a$ lần thứ ba, rồi quay mặt số theo chiều ngược lại cho đến khi điểm mốc gặp vạch số $b$ lần thứ hai, cuối cùng quay mặt số ngược chiều kim đồng hồ cho đến khi điểm mốc gặp vạch số $c$ lần đầu tiên. Nếu $a, b, c$ phải khác nhau đôi một, thì có bao nhiêu cách chọn mật mã cho khoá tổ hợp trên?
	}{ 
	\begin{tikzpicture}[line join=round,line cap=round,>=Stealth,thick,x=1cm,y=1cm,font=\footnotesize]
			\draw[black,thick,fill=gray!20]  (-1,0) -- (-1,1.7)
				arc (180:0:1cm) -- (1,0) --  cycle;
			\draw[black,thick,fill=white]  (-.7,0) -- (-.7,1.7)
				arc (180:0:.7cm)  -- (.7,0) -- cycle;
			\draw[gray!70,fill=gray!10] (0,0) circle (1.4cm); 
			\draw[gray,fill=gray!30] (70:1.1cm) arc (70:-250:1.1cm) -- (-250:1.3cm) arc (110:70:1.3cm)  -- cycle
				(90:1cm)  -- (90:1.2cm);
			\draw[fill=black] (0,0) circle  (1cm);
			\draw[very thin,white,fill=gray!70] (0,0) circle  (5mm);
			\fill[gray] (0,0)  circle (4.5mm);
			\foreach \i  in {0,1,...,40}{
				\draw[white] (\i*9:9mm)  -- (\i*9:1cm);
			} 
			\foreach \i/\j in  {-5/0,-4/5,-3/10,-2/15,-1/20,0/25,1/30,2/35}{
				\node  at (\i*45:7.5mm) [rotate={(\i-2)*45}]{\color{white}\tiny\j};
			 }
		\end{tikzpicture}
	 }
	\loigiai{ 
	 Có 40 cách chọn số $a$ từ các số từ 0 đến 39. Tiếp theo, có 39 cách chọn số $b$ từ 39 số còn lại. Cuối cùng, có 38 cách chọn số $c$ từ 38 số còn lại. Áp dụng quy tắc nhân, ta có 40.39.38 = 59280 cách chọn mật mã cho khoá.}
\end{bt}

% nguồn: MaHoaNH_TuLuan_Lop10
\begin{bt}%[0D8Y1-2]%[ma:41d7187e]
	\immini{ 
	Tung một con xúc xắc ba lần liên tiếp  và ghi lại kết quả (chẳng hạn, $2-5-4$ nếu số chấm xuất hiện lần lượt là 2,5 và 4). Có tất cả bao nhiêu kết quả khác nhau có thể xảy ra?
	 }{
	\begin{tikzpicture}[line join=round,line cap=round,>=Stealth,thick,x=1cm,y=1cm,font=\footnotesize,scale=.4]
			\path 
				(0,0) coordinate (goc) 
				(30:2cm)  coordinate (xs1)
				 (150:2cm) coordinate (xs2)
				(-90:2cm)  coordinate (xs3)
				($(xs1)+(xs2)-(goc)$)  coordinate (nap)
				($(xs1)+(xs3)-(goc)$) coordinate (phai) 
				($(xs2)+(xs3)-(goc)$) coordinate  (trai);
			\draw[rounded corners = 4pt,fill=gray] (xs3)  -- (phai) -- (xs1) -- (nap) -- (xs2) -- (trai) -- cycle;
			\draw[rounded corners = 4pt,fill=gray!10] (goc) -- (xs1) -- (nap)  -- (xs2) -- cycle;
			\draw[rounded  corners = 4pt,fill=gray!70] (goc) -- (xs3) -- (phai) -- (xs1) -- cycle;
			\draw[rounded corners = 4pt,fill=gray!30] (goc) -- (xs3) -- (trai) -- (xs2) --  cycle;
			\foreach \i in {0.33,0.67} 
			\draw[fill=black] ($(xs2)!\i!(xs3)$) circle  (3.5pt);
			\foreach \i in  {0.25,0.5,0.75}
			\draw[fill=black] ($(xs1)!\i!(xs3)$) circle (3.5pt); 
			\foreach \i in {0.33,0.67} 
			\draw[fill=black]  ($(xs2)!\i!(xs1)$) circle (3.5pt)
				($(goc)!\i!(nap)$) circle  (3.5pt);
		\end{tikzpicture}
	} 
	 \loigiai{
	Có thể coi việc tung con xúc xắc ba lần liên tiếp là công việc gồm ba công đoạn, mỗi công đoạn là một lần tung. Mỗi lần tung đều có 6 khả năng khác nhau xảy ra (số chấm xuất hiện là $1;2;\ldots;6$). Do đó, theo quy tắc nhân, ta có $6 \cdot 6 \cdot  6=216$ kết quả khác nhau có thể xảy ra.}
\end{bt}

% nguồn: MaHoaNH_TuLuan_Lop10
\begin{bt}%[0D8T1-2]%[ma:27276564]
	 \immini{
	Mã số nhân viên của một công  ty có 4 kí tự, gồm một chữ cái đầu tiên (từ 6 chữ cái A, B, C, D, E, F) và tiếp theo là 3 chữ số (từ các chữ số 0; 1; \ldots; 9). Công ty có thể tạo ra bao nhiêu mã số nhân viên theo cách này?
	}{ 
	\begin{tikzpicture}[line join=round,line cap=round,>=Stealth,thick,x=1cm,y=1cm,font=\footnotesize]
			 \draw[rounded corners,fill=gray!20] (0,0) rectangle (5,2.5);
			 \draw (0,2) -- (5,2)
				(2.5,2)  -- (2.5,0)
				(3.75,2.1) node[anchor=north]{\tiny  THẺ NHÂN VIÊN};
			\foreach  \i in {1.5,1.2,0.9}{
				 \draw (2.8,\i) -- (4.7,\i);
			 }
			\draw  (3.7,0.6) rectangle (2.8,0.3)
				(3.8,0.6) rectangle (4.7,0.3); 
			\node[rounded  corners,draw=gray] at (1.25,0.3) {\tiny\sffamily\bfseries MS: ABEF239};
			\draw[fill=gray!10]  (0.4,0.6) rectangle (2.1,1.8);
			\draw (.9,1.3) to[out=90,in=150] (1.4,1.6) 
				to[out=30,in=90] (1.6,1.3) 
				(1.45,1.5).. controls  ++(-100:3mm) and ++(120:5mm).. (.9,1.1) coordinate (mt)
				(1.45,1.5).. controls ++(-90:3mm) and  ++(60:4mm).. (1.6,1.1) coordinate (mp)
				(mt).. controls ++(-70:5mm) and ++(-110:5mm).. (mp) 
				coordinate[pos=.4](ct) 
				 coordinate[pos=.6](cp);
			\draw (ct) to[out=-90,in=20] (1,0.6) 
				(cp) to[out=-90,in=160] (1.5,0.6); 
		\end{tikzpicture}
	 }
	\loigiai{ 
		Có 6 cách chọn chữ  cái cho kí tự đầu tiên.\\
		Với  3 kí tự tiếp theo, mỗi kí tự có 10 cách chọn từ 10 chữ số $0,1,2,\ldots,9$.\\
	Theo quy tắc  nhân, công ty có thể tạo ra $6 \cdot 10 \cdot 10 \cdot 10=6000$ mã số nhân viên.}
\end{bt}
